\section{Discussion}
\label{sec:discussion}

The analysis isolates three levels that should not be conflated.
First, \Cref{ass:drift} makes the stochastic prior well defined.
Second, \Cref{ass:girsanov} gives a finite-horizon change of measure and
a terminal density.  Third, posterior contraction requires local mass,
tests, and identifiability for a specified sequence of experiments.
Passing one level does not automatically establish the next.

The lifted approximation in \Cref{thm:discrete} resolves a specific
measure-theoretic obstruction.  It is suited to Hellinger comparison
because both posteriors remain measures on $\HH$.  A practical solver
may store only the first $N$ coordinates; its output should then be
compared after projection or in a topology compatible with that
finite-dimensional representation.  Hellinger convergence to the full
posterior cannot hold for the zero-padded law.

Several questions remain open.  A score-matching theorem would be
needed to relate a particular training objective to the learned drift
and to $\varepsilon_N$.  An unconditional contraction theorem would
need lower control of the terminal Girsanov density near the truth,
together with model-specific tests and sieves.  Nonlinear forward maps
would require separate likelihood bounds and inverse stability
assumptions.  Finally, sampler accuracy is distinct from posterior
approximation and must include time discretization and mixing errors.

The conclusions here are consequently modest but checkable: a stable
posterior for a well-defined learned OU drift prior, a quantitative
full-space discretization bound, and an explicit account of the
additional assumptions needed for statistical recovery.
